\documentclass[12pt]{article}
\usepackage{graphicx} % Required for inserting images
\usepackage[margin=0.5in]{geometry}
\usepackage{amsmath, amssymb}
\usepackage{mathrsfs}


\usepackage{soul}


% Dr. David Brown Commands
\newcommand{\ds}{\displaystyle}
\newcommand{\R}{\mathbb{R}}
\newcommand{\N}{\mathbb{N}}
\newcommand{\Q}{\mathbb{Q}}
\newcommand{\C}{\mathbb{C}}


% Commands to get script r
\usepackage{calligra}
\DeclareMathAlphabet{\mathcalligra}{T1}{calligra}{m}{n}
\DeclareFontShape{T1}{calligra}{m}{n}{<->s*[2.2]callig15}{}
\newcommand{\scripty}[1]{\ensuremath{\mathcalligra{#1}}}

% Unit commands
\newcommand{\degree}{\text{°}}
\newcommand{\unitSpeed}{\frac{\text{m}}{\text{s}}}
\newcommand{\unitAcceleration}{\frac{\text{m}}{\text{s}^2}}

% Constant commands
\newcommand{\avogadro}{6.022\times10^{23}}

\newcommand{\boltzmannConstK}{1.381\times10^{-23}\frac{\text{J}}{\text{K}}}
\newcommand{\boltzmannConstInEV}{8.617\times10^{-5}\frac{\text{eV}}{\text{K}}}
\newcommand{\planckConst}{6.626\times10^{-34}\text{J}\cdot\text{s}}

\newcommand{\atmP}{1.01325\times10^5\,\text{Pa}}

% Commands to easily write partial derivatives
\newcommand{\partDeriv}[2]{\frac{\partial #1}{\partial #2}}
\newcommand{\doublePartDeriv}[2]{\frac{\partial^2 #1}{\partial^2 #2}}


\title{Problem 7.19}
\author{Gabriel Decker}


\begin{document}


\maketitle
\begin{flushleft}



In this problem, we will consider whether the electrons in a chunk of copper can be considered a degenerate fermi gas.
We will calcuate the Fermi energy, Fermi temperature, degeneracy pressure, and the contribution of the degeneracy pressure to the bulk modulus.
We will assume each atom of copper in a chunk of copper contributes one electron.
We will use the following data for copper: density $\rho=8960\,\text{kg/m}^3$ and atomic mass $m_a=0.063546\,\text{kg/mol}$.\\~\\

The relevant formulas are the following.
\begin{equation}
    \epsilon_F=\frac{h^2}{8m}\left(\frac{3N}{\pi V}\right)^{2/3}
\end{equation}
\begin{equation}
    T_F=\frac{\epsilon_F}{k}
\end{equation}
\begin{equation}
    P=\frac{2}{3}\frac{U}{V}
\end{equation}
\begin{equation}
    B=\frac{10}{9}\frac{U}{V}
\end{equation}
where $B$ is the bulk modulus and
\begin{equation}
    U=\frac{3}{5}N\epsilon_F
\end{equation}\\~\\

We aren't given a specific amount of copper to deal with, but that doesn't matter.
What matters is the number of electrons contained in a volume.
So, let's calculate $N/V$ using the atomic mass and density because that will be very important throughout.
We can use atomic mass because we know that each atom will contribute one electron to the cloud.
Recall that $\rho=\frac{m}{V}$ and $m_a$ is the mass per $\avogadro$ particles.
$$\rho=\frac{m}{V}$$
$$\implies\frac{N}{V}=\frac{N\rho}{m}$$
If we consider $N=\avogadro=1\,\text{mol}$, then this implies that $m=m_a\cdot1\,\text{mol}$. Therefore, we get the following.
$$\implies\frac{N}{V}=\frac{(1\,\text{mol})(8960\,\text{kg/m}^3)}{(0.063546\,\text{kg/mol})(1\,\text{mol})}$$
$$\implies\frac{N}{V}=\frac{(\avogadro\,\text{electrons})(8960\,\text{kg/m}^3)}{(0.063546\,\text{kg/mol})(1\,\text{mol})}$$
$$\implies\frac{N}{V}=8.49103\times10^{28}\,\text{electrons/m}^3$$\\~\\

This ratio will be used to find the quantities we're actually interested in.
First, let's find the Fermi energy.
Note that from now on, $m$ refers to the mass of an electron $m=9.1094\times10^{-31}\,\text{kg}$.
$$\epsilon_F=\frac{h^2}{8m}\left(\frac{3N}{\pi V}\right)^{2/3}$$
$$\implies\epsilon_F=\frac{(\planckConst)^2}{8\cdot(9.1094\times10^{-31}\,\text{kg})}\left(\frac{3}{\pi}(8.49103\times10^{28}\,\text{1/m}^3)\right)^{2/3}$$
$$\implies\epsilon_F=1.1286\times10^{-18}\,\text{J}$$\\~\\

Now for the Fermi Temperature.
$$T_F=\frac{\epsilon_F}{k}$$
$$\implies T_F=\frac{1.1286\times10^{-18}\,\text{J}}{\boltzmannConstK}$$
$$\implies T_F=81724\,\text{K}$$
Based on this, room temperature is certainly sufficiently low to treat the system as a degenerate electron gas.\\~\\

Now for the degeneracy pressure.
First, let's simplify our calculation a bit by combining our definition of $U$ into $P$.
$$P=\frac{2}{3}\frac{U}{V}$$
$$\implies P=\frac{2}{5}\epsilon_F\cdot\frac{N}{V}$$
$$\implies P=\frac{2}{5}(1.1286\times10^{-18}\,\text{J})\cdot(8.49103\times10^{28}\,\text{1/m}^3)$$
$$\implies P=3.833\times10^{10}\,\text{Pa}$$\\~\\

And for the bulk modulus, we'll do the following since it's so similar to $P$.
$$B=\frac{10}{9}\frac{U}{V}$$
$$\implies B=\frac{10}{9}\frac{3}{2}\cdot\frac{2}{3}\frac{U}{V}$$
$$\implies B=\frac{5}{3}P$$
$$\implies B=\frac{5}{3}(3.833\times10^{10}\,\text{Pa})$$
$$\implies B=6.389\,\text{Pa}$$






\end{flushleft}

\end{document}
